% Inicio del preámbulo

\documentclass[a4paper,11pt]{article} %Modifica el tipo de documento y el tamaño de la letra.
\usepackage{fontspec}
\usepackage{setspace}
\setmainfont{Times New Roman}
\usepackage{titlesec}
%\setmainfont{Arial}
\usepackage[utf8]{inputenc} %Formato UTF-8 para caracteres especiales.
\singlespacing
\usepackage{lmodern}
\usepackage[
  style = ieee,
  backend = biber,
  url=true,  
  isbn = false,
  maxbibnames = 6, % Número máximo de autores que pueden aparecer en una referencia (6)
  minnames=1,
  citestyle=numeric] % Número mínimo de autores que aparecerán al truncar: autor1 et al.
{biblatex}

\usepackage[dvipsnames]{xcolor}
\usepackage[spanish,es-tabla]{babel}
\usepackage{listings}

\usepackage[a4paper, left=2cm, right=2cm, top=2.5cm, bottom=2.5cm]{geometry}

\usepackage[shortlabels]{enumitem}
\usepackage{amsmath,amssymb,amsfonts,latexsym,cancel}
\usepackage{hyperref}
\usepackage{subcaption}
\usepackage{wrapfig}
\usepackage[rflt]{floatflt}
%\usepackage[pdftex]{graphicx}
\usepackage{graphicx}
\usepackage{fancyhdr} %Paquete para el header y el formato de la portada. No sugiero borrarlo!
\usepackage{float}
\usepackage{longtable,multirow,booktabs}
\usepackage[dvipsnames,table]{xcolor}
\usepackage{tabularx}
\usepackage{caption}
\usepackage{sidecap}
\usepackage{adjustbox}
\usepackage{parskip}
\usepackage{enumitem}
\usepackage{tikz}
\usetikzlibrary{arrows.meta, positioning, shapes.geometric, shapes, calc}
\usepackage{lipsum}
\usepackage{wrapfig} %texto alrededor de Images \begin{wrapfigure}{alignment}{width}
\usepackage{xlop}	%divisiones con caja \opidiv{316}{50}

\usepackage{multicol} %varias columnas


\usepackage{textcomp}   %poner símbolo copyright

\usepackage{endnotes}


% Para hacer recuadros
% https://ondahostil.wordpress.com/2017/05/17/lo-que-he-aprendido-cuadros-de-texto-de-colores-en-latex/
\usepackage{tcolorbox}

%%% Matemáticas Paquetes AMS
\usepackage{amsfonts}
\usepackage{amsmath} % Matemáticas
\usepackage{amssymb,amsthm} % Símbolos y teoremas
\usepackage{mathtools} % % Algunos añadidos y correcciones a amsmath
\usepackage{upgreek} % Letras griegas
\usepackage{esvect}
\usepackage{dsfont} % Conjunto unario
\usepackage{multirow} % Tablas con casillas de varias alturas

%Fin de Préambulo